\documentclass[11pt]{article}
\usepackage{graphicx}
\usepackage{amsmath, amssymb}
\usepackage{hyperref}
\usepackage{enumitem}
\usepackage{geometry}
\geometry{margin=1in}

\title{Reconfigurable Computing --- Post Mid-Semester Lecture Notes (Lectures 19--26)}
\author{Based on course slides by Dr.~Meetha V.~Shenoy}
\date{}

\begin{document}
\maketitle
\tableofcontents
\newpage

%%%%%%%%%%%%%%%%%%%%%%%%%%%%%%%%%%%%%%%%%%%%%%%%%%%%%%%%%%%%%
\section{Lecture 19: FPGA Routing and Architecture}
%%%%%%%%%%%%%%%%%%%%%%%%%%%%%%%%%%%%%%%%%%%%%%%%%%%%%%%%%%%%%

\subsection{Overview of FPGA CAD Flow}
According to the flow diagram shown in the slides
(Design Entry $\rightarrow$ Synthesis $\rightarrow$ Logic Optimization $\rightarrow$ Mapping to k-LUT
$\rightarrow$ Packing LUTs to CLBs $\rightarrow$ Placement $\rightarrow$ Routing $\rightarrow$ Simulation
$\rightarrow$ Configuration)
:contentReference[oaicite:1]{index=1}, the FPGA compilation process proceeds through the following steps:

\begin{enumerate}
    \item \textbf{Design Entry}: HDL / schematic input.
    \item \textbf{Synthesis}: Converts HDL to gate-level netlist.
    \item \textbf{Logic Optimization}.
    \item \textbf{Technology Mapping}: Map to k-input LUTs.
    \item \textbf{Packing}: Pack LUTs into CLBs.
    \item \textbf{Placement}: Assign CLBs to FPGA grid.
    \item \textbf{Routing}: Connect all placed blocks using routing fabric.
    \item \textbf{Bitstream Generation}.
\end{enumerate}

\subsection{Island-Style FPGA Architecture}
The slides depict an island-style FPGA (logic blocks surrounded by routing channels and switch boxes)
:contentReference[oaicite:2]{index=2}.

\begin{itemize}
    \item \textbf{Logic Blocks (L)} implement combinational/sequential logic.
    \item \textbf{Connection Boxes (C)} connect CLB pins to routing tracks.
    \item \textbf{Switch Boxes (S)} connect routing tracks to each other.
\end{itemize}

\subsection{Connection Block Flexibility}
The connection block has a parameter:
\[
F_c = \text{number of routing tracks each CLB pin may connect to}
\]
Higher $F_c$ increases routability but increases area. :contentReference[oaicite:3]{index=3}

\subsection{Switch Box Types}
Switch boxes connect horizontal and vertical tracks. Two major types shown in the slides are
:contentReference[oaicite:4]{index=4}:
\begin{itemize}
    \item \textbf{Planar / Subset switch box}: Only maps track $i$ to track $i$ (0--0, 1--1).
    \item \textbf{Wilton switch box}: Permits permutation across domains (0--3, 1--2, etc.).
\end{itemize}

\subsection{Routing as a Grid Graph}
Slides describe routing as a graph where:
\begin{itemize}
    \item Nodes = switches, pins, wire segments.
    \item Obstacles = used routing resources.
    \item Cost = wirelength + congestion + timing. :contentReference[oaicite:5]{index=5}
\end{itemize}

\subsection{Maze Routing Algorithm}
A BFS/Dijkstra style expansion is used:
\begin{enumerate}
    \item Initialize grid.
    \item Expand wave from source.
    \item When destination reached, backtrack.
    \item Optimize for congestion/timing. :contentReference[oaicite:6]{index=6}
\end{enumerate}

\subsection{Global vs Detailed Routing}
\begin{itemize}
    \item \textbf{Global routing}: Choose coarse path (which channels).
    \item \textbf{Detailed routing}: Assign specific tracks and switches. :contentReference[oaicite:7]{index=7}
\end{itemize}

\subsection{Pathfinder Algorithm}
Iterative negotiated routing:
\begin{enumerate}
    \item Initial shortest‐path routing.
    \item Compute congestion.
    \item Increase cost of overused wires.
    \item Rip-up and reroute.
    \item Iterate until convergence. :contentReference[oaicite:8]{index=8}
\end{enumerate}

\newpage

%%%%%%%%%%%%%%%%%%%%%%%%%%%%%%%%%%%%%%%%%%%%%%%%%%%%%%%%%%%%%
\section{Lecture 20: ZedBoard, Zynq SoC and JTAG}
%%%%%%%%%%%%%%%%%%%%%%%%%%%%%%%%%%%%%%%%%%%%%%%%%%%%%%%%%%%%%

\subsection{ZedBoard Overview}
Slides provide a labeled diagram of the ZedBoard main components
(USB-JTAG, USB-UART, GPIO, HDMI, Ethernet, Memory interfaces)
:contentReference[oaicite:9]{index=9}.

\subsection{Zynq-7000 Architecture}
Zynq integrates:
\begin{itemize}
    \item \textbf{Processing System (PS)}: Dual-core ARM Cortex-A9.
    \item \textbf{Programmable Logic (PL)}: Artix-7 FPGA fabric.
    \item \textbf{AXI Interconnect}: Communication between PS and PL. :contentReference[oaicite:10]{index=10}
\end{itemize}

\subsection{JTAG and IEEE 1149.1}
The wrapper architecture figure shows boundary scan cells around the internal logic
:contentReference[oaicite:11]{index=11}.

\subsubsection*{Mandatory pins}
\begin{itemize}
    \item TCK, TMS, TDI, TDO, optional TRST. :contentReference[oaicite:12]{index=12}
\end{itemize}

\subsection{JTAG Instructions}
\begin{itemize}
    \item \textbf{EXTEST}: Tests external connections (figure in slides).
    \item \textbf{BYPASS}: 1-bit bypass register.
    \item \textbf{INTEST}: Test internal logic. :contentReference[oaicite:13]{index=13}
\end{itemize}

\subsection{Programming Methods}
ZedBoard supports programming via:
\begin{enumerate}
    \item USB-JTAG (default).
    \item JTAG header.
    \item SPI Flash.
    \item SD Card. :contentReference[oaicite:14]{index=14}
\end{enumerate}

\newpage

%%%%%%%%%%%%%%%%%%%%%%%%%%%%%%%%%%%%%%%%%%%%%%%%%%%%%%%%%%%%%
\section{Lecture 21: Security of FPGA-Based Systems (Part 1)}
%%%%%%%%%%%%%%%%%%%%%%%%%%%%%%%%%%%%%%%%%%%%%%%%%%%%%%%%%%%%%

\subsection{Supply Chain Model}
Slides show interactions between foundries, FPGA vendors, system developers, IP vendors, end users
:contentReference[oaicite:15]{index=15}.

\subsection{Threat: Cloning}
FPGAs load bitstreams from external flash; attacker can read this bitstream and reprogram unauthorized devices
:contentReference[oaicite:16]{index=16}.

\subsection{Identification: Friend or Foe}
Authentication protocol:
\begin{enumerate}
    \item FPGA sends challenge $C$ to EEPROM.
    \item EEPROM computes $\mathrm{Hash}(K_{\text{EEPROM}}, C)$.
    \item FPGA computes $\mathrm{Hash}(K_{\text{FPGA}}, C)$.
    \item If match $\Rightarrow$ valid device. :contentReference[oaicite:17]{index=17}
\end{enumerate}

\subsection{Physically Unclonable Functions (PUFs)}
Key points from the slides:
\begin{itemize}
    \item Ring oscillators differ due to manufacturing variations.
    \item Frequency differences yield unique response bits.
    \item Multiple challenges produce a device fingerprint. :contentReference[oaicite:18]{index=18}
\end{itemize}

\subsection{Hardware Trojans}
A malicious modification consisting of:
\begin{itemize}
    \item \textbf{Trigger}: Rare conditions.
    \item \textbf{Payload}: Key leakage, disabling chip, reduced performance. :contentReference[oaicite:19]{index=19}
\end{itemize}

\subsection{Backdoor Communication}
Unauthorized hidden channels inside FPGA IP blocks. :contentReference[oaicite:20]{index=20}

\newpage

%%%%%%%%%%%%%%%%%%%%%%%%%%%%%%%%%%%%%%%%%%%%%%%%%%%%%%%%%%%%%
\section{Lecture 22: Security of FPGA-Based Systems (Part 2)}
%%%%%%%%%%%%%%%%%%%%%%%%%%%%%%%%%%%%%%%%%%%%%%%%%%%%%%%%%%%%%

\subsection{Replay Attacks}
Attacker replays an older bitstream version to exploit vulnerabilities
:contentReference[oaicite:21]{index=21}.

PUF-based license generation:
\[
\text{License}_1 = f(R, \text{Version}, \text{Secret})
\]
Bitstream executes only when PUF-derived license matches. :contentReference[oaicite:22]{index=22}

\subsection{Secure Boot}
Slides illustrate standard secure-boot flow:
\begin{itemize}
    \item Host signs image.
    \item Device verifies signature using public key in fuses. :contentReference[oaicite:23]{index=23}
\end{itemize}

\subsection{Reverse Engineering of ASICs}
Uses delayering, SEM imaging, and X-ray tomography to restore netlists. :contentReference[oaicite:24]{index=24}

\subsection{IC Camouflage}
Gates that appear identical physically but behave differently. :contentReference[oaicite:25]{index=25}

\subsection{Logic Obfuscation}
Insert key-dependent gates; wrong keys produce incorrect behavior. :contentReference[oaicite:26]{index=26}

\subsection{Tamper Detection}
Secure FPGAs include:
\begin{itemize}
    \item Light sensors.
    \item Switches.
    \item Vibration detectors.
    \item Zeroization of secrets when triggered. :contentReference[oaicite:27]{index=27}
\end{itemize}

\newpage

%%%%%%%%%%%%%%%%%%%%%%%%%%%%%%%%%%%%%%%%%%%%%%%%%%%%%%%%%%%%%
\section{Lecture 23: Coarse-Grained Reconfigurable Architectures (CGRAs)}
%%%%%%%%%%%%%%%%%%%%%%%%%%%%%%%%%%%%%%%%%%%%%%%%%%%%%%%%%%%%%

\subsection{CGRA Structure}
Slides depict Processing Element (PE) arrays, register files, configuration memories, host processor
:contentReference[oaicite:28]{index=28}.

\subsection{Processing Elements}
Each PE contains:
\begin{itemize}
    \item ALU + shifter.
    \item Local register file.
    \item Shared register for neighbor communication. :contentReference[oaicite:29]{index=29}
\end{itemize}

\subsection{Execution Controller}
Controls:
\begin{itemize}
    \item PE activation schedule.
    \item Context switching.
    \item Pipeline start/stop.
    \item Synchronization. :contentReference[oaicite:30]{index=30}
\end{itemize}

\subsection{CGRAs vs FPGAs}
\begin{itemize}
    \item FPGA: logic-level programming.
    \item CGRA: operation-level programming (dataflow-like). :contentReference[oaicite:31]{index=31}
\end{itemize}

\subsection{Interconnect Types}
\begin{enumerate}
    \item \textbf{Nearest Neighbor}.  
    \item \textbf{Next-Hop / Flexible Switchbox}.  
    \item \textbf{Bus-Based Interconnect}.  
    \item \textbf{Augmented Interconnect}. :contentReference[oaicite:32]{index=32}
\end{enumerate}

\newpage

%%%%%%%%%%%%%%%%%%%%%%%%%%%%%%%%%%%%%%%%%%%%%%%%%%%%%%%%%%%%%
\section{Lecture 24--26: Configurable Logic Block (CLB) Architecture}
%%%%%%%%%%%%%%%%%%%%%%%%%%%%%%%%%%%%%%%%%%%%%%%%%%%%%%%%%%%%%

\subsection{CLB Structure}
Slides show:
\begin{itemize}
    \item 1 CLB = 2 slices.
    \item Each slice = 4 LUTs + 8 FFs + carry chain + multiplexers. :contentReference[oaicite:33]{index=33}
\end{itemize}

There are two slice types:
\begin{itemize}
    \item \textbf{SLICEM}: supports distributed RAM + SRL.
    \item \textbf{SLICEL}: logic only. :contentReference[oaicite:34]{index=34}
\end{itemize}

\subsection{LUT Functionality}
\begin{itemize}
    \item 6-input LUT with output $O_6$.
    \item Can act as two 5-input LUTs using $O_5$. :contentReference[oaicite:35]{index=35}
\end{itemize}

\subsection{LUTs as Multiplexers}
Slides illustrate 4:1, 8:1, and 16:1 multiplexer construction using LUTs  
:contentReference[oaicite:36]{index=36}.

\subsection{Storage Elements}
\begin{itemize}
    \item Each LUT output feeds optional FF.
    \item Total: 8 FFs per slice. :contentReference[oaicite:37]{index=37}
\end{itemize}

\subsection{Distributed RAM}
Available only in SLICEM. Good for:
\begin{itemize}
    \item Small tables.
    \item Coefficients.
    \item Low-latency lookup. :contentReference[oaicite:38]{index=38}
\end{itemize}

\subsection{Shift Register LUTs (SRLs)}
\subsubsection*{SRL32}
A LUT can become a 32-bit shift register:
\begin{itemize}
    \item Depth = 1–32.
    \item Address lines A[4:0] select tap.
    \item Q31 provides cascading. :contentReference[oaicite:39]{index=39}
\end{itemize}

\subsubsection*{SRL Cascading}
Up to 128 cycles using 4 LUTs (SRL1--SRL4). Tables and diagrams shown in slides illustrate ranges
:contentReference[oaicite:40]{index=40}.

\subsection{Block RAM vs Distributed RAM}
\begin{center}
\begin{tabular}{|c|c|c|}
\hline
Characteristic & Distributed RAM & Block RAM \\
\hline
Depth & $\le 128$ & large (512+) \\
Width & $\le 8$ bits & up to 72 bits \\
Read & async & sync \\
Location & LUT fabric & BRAM tile \\
\hline
\end{tabular}
\end{center}
:contentReference[oaicite:41]{index=41}

\subsection{Pipelining}
Breaking long combinational paths into stages improves:
\begin{itemize}
    \item Frequency.
    \item Throughput.
    \item Timing closure. :contentReference[oaicite:42]{index=42}
\end{itemize}

Example:
\[
r = (a+b)\cdot c \quad\Rightarrow\quad
\begin{aligned}
    \text{Stage 1:}&\ a+b \\
    \text{Stage 2:}&\ (a+b)\cdot c
\end{aligned}
\]

\newpage

%%%%%%%%%%%%%%%%%%%%%%%%%%%%%%%%%%%%%%%%%%%%%%%%%%%%%%%%%%%%%
\section{Lecture 26 Addendum: FPGA Implementation of 3×3 Convolution}
%%%%%%%%%%%%%%%%%%%%%%%%%%%%%%%%%%%%%%%%%%%%%%%%%%%%%%%%%%%%%

\subsection{Problem Definition}
Compute 3×3 convolution over 28×28 image:
\begin{itemize}
    \item 9 multiplies + 8 adds/pixel. :contentReference[oaicite:43]{index=43}
\end{itemize}

\subsection{Resource Usage}
\begin{itemize}
    \item \textbf{Line buffers}: BRAM.
    \item \textbf{Kernel coefficients}: Distributed RAM / registers.
    \item \textbf{Window shift}: SRL32.
    \item \textbf{Multiply}: DSP48E1.
    \item \textbf{Adder tree}: CARRY4 chain. :contentReference[oaicite:44]{index=44}
\end{itemize}

\subsection{Pipeline Stages}
\begin{enumerate}
    \item Window extraction (BRAM/RAM).
    \item 9 parallel multiplications.
    \item Row partial sums.
    \item Final accumulation.
    \item Activation/normalization.
    \item Output store. :contentReference[oaicite:45]{index=45}
\end{enumerate}

\end{document}

